\chapter{Lone Working, Site Security and Violence at Work Risk Assessment}
\label{chap:Lone Working, Site Security and Violence at Work Risk Assessment}

%---------------------------- SECTION ----------------------------
\section{Why this assessment matters in construction}

\paragraph{Lone working in construction is far more prevalent than is commonly acknowledged. The surveyor conducting a condition survey on an unoccupied commercial property, the site manager making an early-morning walk around before the workforce arrives, the plumber completing a final connection in a plant room while the rest of the site has gone home, and the security guard conducting an overnight patrol of a partially completed residential block are all lone workers. The fundamental characteristic of lone working --- the absence of a colleague who can summon help or provide assistance if the worker is injured, taken ill, confronted by an intruder, or trapped --- creates a specific risk profile that must be assessed separately from the general site risk assessment.}

\paragraph{Violence at work in construction takes two distinct forms: violence from members of the public, intruders, or third parties at or adjacent to the construction site; and violence from colleagues, supervisors, or subcontractors within the construction workforce. The construction environment presents specific violence risk factors: urban construction sites in public spaces attract the attention of individuals in crisis, intoxicated individuals, and organised criminal networks who target valuable plant, fuel, and materials; sites under construction in socially deprived areas may attract hostility from local residents who are affected by noise, dust, and disruption; and the hierarchical and pressure-laden culture of construction can create conditions in which aggressive management behaviour or operative-to-operative conflict escalates to physical confrontation.}

\paragraph{Site security is the physical and procedural framework that protects the construction site, its workers, and its assets from intruder access, theft, vandalism, and the particular hazard of members of the public accessing an inherently dangerous environment. The Construction Industry Research and Information Association (CIRIA) estimates that theft from UK construction sites costs the industry several hundred million pounds annually; beyond the financial cost, the security failure that allows an intruder to access an active construction site creates a serious risk of injury or death --- to the intruder who falls through an unprotected floor opening, contacts live services, or is struck by plant, and to the site operatives who may encounter the intruder in a hostile exchange.}

\paragraph{Hostile Vehicle Mitigation (HVM) has become an increasing feature of construction site security in urban environments following several high-profile incidents in which vehicles have been used as weapons in public spaces. A construction site in a city centre street, with its perimeter hoarding at the footpath edge, its delivery vehicle access points, and its proximity to pedestrianised areas, presents a genuine HVM challenge. The temporary nature of construction site boundaries and the operational requirement for large vehicles to enter and exit the site creates specific vulnerabilities that permanent building security arrangements do not face. The security risk assessment must address HVM as a distinct risk pathway, not merely as an element of general hoarding provision.}

\barquote{``A lone worker in difficulty is invisible to the world until someone checks on them --- or until they are found. The lone worker risk assessment must ensure that the interval between those two possibilities is measured in minutes, not hours.''}

\example{Lone Working Scenario}{A mechanical and electrical contractor's electrician is conducting cable testing in a basement plant room of a recently completed office building on a Saturday morning, when the main contractor's site offices are unstaffed. She is working alone and has not followed the company's check-in procedure because the site manager informally told her ``the job shouldn't take more than an hour.'' After ninety minutes she slips on a wet floor, fractures her wrist, and is unable to climb the ladder from the basement plant room. Her mobile phone has no signal in the basement. She is not found for four hours, when the security guard conducting an end-of-shift patrol notices her vehicle still in the car park. She has developed early signs of hypothermia. The investigation identifies: no lone worker risk assessment for the specific activity; no check-in procedure; no phone signal assessment for the basement area; and no agreement with the security provider for welfare checks on lone workers.}

%---------------------------- SECTION ----------------------------
\section{Legal and professional framework}

\paragraph{The primary legal basis for lone worker risk management is the Management of Health and Safety at Work Regulations 1999, which require the employer to assess the health and safety risks to lone workers, consider whether the risks can be adequately controlled, and put in place appropriate monitoring and communication systems. There is no specific lone worker regulation; the risk assessment and control duty under the Management Regulations, and the general employer duty under the Health and Safety at Work etc.\ Act 1974, together impose the obligation to manage lone worker risk as part of the broader duty of care.}

\paragraph{Violence at work is addressed under the same general framework: the employer's duty under Section 2 of the Health and Safety at Work etc.\ Act 1974 includes the duty to protect employees from violence that is a foreseeable risk of their work. Where violence involves an assault, the reporting obligations under RIDDOR 2013 apply. The Reporting of Injuries, Diseases and Dangerous Occurrences Regulations 2013 require reporting of any accident at work which results in a worker being away from work or unable to perform their normal duties for more than seven days (over-seven-day injury), and any fatal injury.}

\begin{itemize}
  \bitem{Management of Health and Safety at Work Regulations 1999}{Regulation 3 requires a suitable and sufficient risk assessment of the risks to lone workers. The assessment must consider the specific hazards faced by the lone worker, the adequacy of any emergency response given the lone nature of the work, and whether the work can safely be done by one person. Some work categorically cannot be done safely by a lone worker (confined space entry, work on live high-voltage systems, work in certain atmospheres) and the risk assessment must identify these categories and impose a two-person minimum for them.}
  \bitem{Health and Safety at Work etc.\ Act 1974}{Section 2 imposes the general duty on the employer to ensure the health and safety of lone workers. The foreseeability of harm to a lone worker from isolation, inability to summon help in an emergency, or violence from third parties is well-established; an employer who sends a worker to conduct lone working without a risk assessment or monitoring system has breached this duty.}
  \bitem{BS 8484:2022 --- Lone Worker Device Standard}{The British Standard for the provision of lone worker device services, specifying the requirements for devices, monitoring centres, and response procedures. Compliance with BS 8484 enables the monitoring centre to request police attendance at grade 1 response (emergency) level without the 999 call process, which significantly reduces response times in genuine lone worker emergencies. The standard covers app-based solutions as well as dedicated devices.}
  \bitem{RIDDOR 2013}{Requires the reporting of violence-related injuries (where the violence arises from the victim's work) in the same way as physical accident injuries. Deaths, specified injuries, and over-seven-day injuries arising from workplace violence must be reported to the HSE. Near-misses involving violence or threatening behaviour should be recorded in the company's incident log even where RIDDOR reporting is not triggered.}
  \bitem{Secure by Design}{The Association of Chief Police Officers (ACPO) and Police Digital Security Centre initiative that applies crime prevention through environmental design principles to construction sites and temporary structures. Secure by Design for construction sites provides guidance on fencing standards, lighting, CCTV, access control, and asset protection that, when applied, significantly reduces the risk of site crime and unauthorised access.}
  \bitem{Counter Terrorism Policing --- Hostile Vehicle Mitigation guidance}{Counter Terrorism Policing provides HVM guidance for construction sites in city-centre and high-footfall locations. The guidance addresses the placement of concrete blocks, water-filled barriers, bollards, and vehicle-activated barriers at site access and egress points to prevent vehicle-borne attack. Where a construction site occupies a significant portion of a public area, the HVM risk assessment should be conducted in consultation with the local Counter Terrorism Security Adviser.}
  \bitem{Protection of Freedoms Act 2012 and UK GDPR 2018}{Govern the use of CCTV on construction sites. CCTV must be used proportionately, must cover only areas where there is a genuine security justification, must be clearly signposted, and must comply with data retention and access requirements. Construction companies must register with the Information Commissioner's Office as data controllers for their CCTV systems.}
  \bitem{CDM 2015}{Requires the Construction Phase Plan to address site security arrangements, including hoarding and fencing standards, access control, and protection of the public from construction hazards. The principal contractor's duty to ensure the site is secure is inherent in the duty to manage the construction site safely.}
\end{itemize}

%---------------------------- SECTION ----------------------------
\section{Conducting the assessment using the five-step method}

\subsection{Step 1 --- Identify Hazards}
\paragraph{Lone working hazards include: sudden illness or physical collapse without a colleague to summon help; injury from a fall, moving plant, structural failure, or manual handling incident with no assistance available; entrapment in a basement, plant room, shaft, or other location with restricted access or no phone signal; assault or threatening behaviour from intruders, members of the public, or hostile third parties; and exposure to hazardous atmospheres (gas, chemical, oxygen deficiency) without a second person to monitor and respond. Security hazards include: unauthorised entry by intruders seeking to steal plant, fuel, copper cable, or catalytic converters from plant; vandalism and arson, including the deliberate ignition of flammable materials on the site; entry by vulnerable persons (children, intoxicated individuals) who are at serious risk of injury in the construction environment; and hostile vehicle attack at site entrances and hoarding boundaries. Violence at work hazards include: physical assault or threatening behaviour by site visitors, members of the public, or adjacent residents; operative-to-operative violence in the context of work disputes, gang-related conflict, or substance-related aggression; and threats made by dissatisfied clients, neighbours, or members of the public to site staff.}

\subsection{Step 2 --- Decide Who or What Is Affected}
\paragraph{Persons affected by lone working risk include any operative or manager who works alone on the site at any time, including before and after normal working hours. Site security personnel are particularly at risk of violence from intruders and should be specifically assessed. Survey personnel, service inspectors, and specialist subcontractors who attend the site individually outside normal working hours are a frequently overlooked lone worker group. Members of the public are at risk from intruder access to the site; children in particular are at risk of fatal injury from falls through unprotected openings, contact with live services, and instability of partially completed structures. The site's immediate neighbours are at risk from the security failure that allows vehicle or structural hazards to extend beyond the site boundary. Plant and equipment assets stored on site overnight are at risk from theft, increasing project costs and causing programme delay.}

\subsection{Step 3 --- Evaluate Likelihood and Consequence}
\paragraph{The consequence of a lone worker emergency --- an operative injured, collapsed, or trapped with no colleague to summon help --- is potentially fatal if the response is delayed beyond the critical window for the medical condition involved (minutes for cardiac arrest, an hour or more for severe blood loss, hours for hypothermia or entrapment without injury). The likelihood of a lone worker incident occurring on a large construction programme over its full duration is significant; the risk assessment must not discount likelihood on the basis that lone working events are rare but must consider the cumulative exposure across the workforce. The consequence of an intruder or public access incident can range from property theft through to the death of a child who accesses the site through inadequate fencing; the latter event is a criminal matter as well as a civil liability, and has resulted in convictions of construction companies and directors.}

\subsection{Step 4 --- Select and Implement Controls}
\paragraph{Controls for lone working must include: a pre-task lone worker risk assessment that confirms the activity can be done safely alone; a check-in system that requires the lone worker to confirm their safety at agreed intervals with a designated contact; a documented escalation procedure specifying what the contact should do if a check-in is missed; the provision of a lone worker device or BS 8484-compliant app that enables the lone worker to raise a silent alarm directly to a monitoring centre; assessment of phone signal coverage at the work location and provision of alternative communication where coverage is absent; and specific restrictions on lone working in confined spaces, at height, on live electrical systems, and in any other category where assistance is immediately required for emergency response. Controls for site security must include: a robust hoarding and fencing installation that meets BS 1722 chain-link or equivalent standards for the site's risk profile; a locking and access control procedure for all site entry and exit points including after-hours access by security personnel; CCTV coverage of the site perimeter, access points, and high-value asset storage areas; a site security plan that is reviewed whenever the site configuration changes; and an out-of-hours intruder response procedure. Controls for violence at work must include: a site behaviour policy that specifically addresses threatening and intimidatory behaviour; a reporting system for violence incidents and near-misses; a personal safety procedure for site staff working in areas of known public hostility; and access to support for workers who have experienced violence at work.}

\subsection{Step 5 --- Record and Review}
\paragraph{The lone worker risk assessment must be reviewed whenever the lone working activity changes (new location, new task, change of working hours), when a lone worker incident occurs, and at least annually as part of the general risk assessment review cycle. The check-in records must be retained as evidence that the monitoring system was operating. The security risk assessment must be reviewed at each significant change to the site configuration (hoarding extension, opening of new access gates, introduction of new high-value assets). CCTV maintenance records and access control logs must be maintained to demonstrate the ongoing functionality of the security controls.}

\example{Five-Step Application}{A principal contractor on an urban residential project operates from 7\,am to 6\,pm with a site manager who arrives at 6\,am and leaves last at 7\,pm. Step 1 identifies: site manager lone working for one hour before and after normal hours; overnight security guard patrolling a site with confirmed copper cable theft in the area; CCTV blind spots at the rear of the site adjacent to an alley. Step 2 identifies: site manager, security guard, and any subcontractor arriving before 7\,am are lone workers. Step 3 rates risk to site manager from medical emergency as High (lone for one hour); risk of violence to security guard as High given recent thefts. Step 4 implements: site manager registers on BS 8484 app before 6\,am solo attendance and checks in at 6\,am and again at 6\,30\,am; security guard equipped with BS 8484 device with 30-minute check-in; CCTV extended to cover rear alley blind spot; concrete barrier blocks installed at front hoarding access point for HVM. Step 5 specifies: monthly security review; immediate review on any intruder incident; annual lone worker policy review.}

%---------------------------- SECTION ----------------------------
\section{Applying the hierarchy of controls}

\paragraph{The hierarchy of controls for lone working, security, and violence must prioritise the elimination or avoidance of the hazardous situation wherever possible. Not all lone working is necessary; task planning that ensures workers are accompanied wherever the risk warrants it is a higher-order control than equipping them with a monitoring device.}

\begin{itemize}
  \bitem{Elimination}{Eliminate lone working where the hazard profile is incompatible with safe lone operation: prohibit lone working in confined spaces absolutely; prohibit lone working on live electrical systems above 50\,V\,AC; require two-person attendance for any early-morning or out-of-hours site access. Eliminate the attractiveness of the site to thieves by removing high-value materials from site overnight, fitting wheel clamps to plant, and removing fuel from mobile plant.}
  \bitem{Substitution}{Substitute extended lone working periods with coordinated working patterns: schedule early-morning and late-evening site tasks so that two workers attend together; substitute unsupervised subcontractor visits with accompanied site access; substitute physical keys for electronically audited access control systems that record entry and exit times.}
  \bitem{Engineering controls}{Install BS 8484-compliant lone worker monitoring devices or apps for all workers who work alone; install CCTV covering the full site perimeter and all access points; install perimeter fencing to an appropriate standard for the site risk profile; install concrete or water-filled HVM barriers at all vehicle access points; install anti-climb paint and defensive landscaping on perimeter fencing; install adequate lighting of site perimeter, access routes, and storage areas for overnight security patrols.}
  \bitem{Administrative controls}{Implement a lone worker register and check-in procedure; produce a lone worker risk assessment for each category of lone working activity; implement a site security plan; conduct site security inductions; maintain a violence and aggression incident log; establish a post-incident support pathway for workers who experience violence; implement a visitor management system that records all site attendees including out-of-hours access; brief all lone workers on the BS 8484 device escalation procedure.}
  \bitem{PPE}{Personal alarm devices (including BS 8484-compliant apps) for lone workers; high-visibility clothing for night-time security patrols; torches and communication radios for security personnel; body-worn cameras for security guards on sites with a history of violent incidents, both as a deterrent and as evidence in criminal proceedings.}
\end{itemize}

%---------------------------- SECTION ----------------------------
\section{Cadence of assessment and review triggers}

\paragraph{Security and lone worker risks on a construction site are dynamic: the site configuration changes as structures are erected, access points are opened and closed, and the value of assets on site varies through the construction programme. The risk assessment cadence must reflect this dynamism and must include specific review points at each phase transition.}

\begin{itemize}
  \bitem{At each phase transition}{The security risk assessment must be reviewed whenever the site moves from one phase to another (groundworks to structure, structure to fit-out, fit-out to handover) as the physical layout, access routes, valuable assets, and lone worker activity patterns all change with each phase.}
  \bitem{After any security incident}{Any intruder access, theft, vandalism, or violence incident must immediately trigger a review of the security risk assessment and the adequacy of the controls. The CCTV footage must be reviewed and any identified blind spots addressed.}
  \bitem{After any lone worker incident}{Any incident involving a lone worker --- injury, entrapment, assault, or a missed check-in that required an emergency response --- must trigger a review of the lone worker risk assessment and a debrief with the affected worker.}
  \bitem{When new high-value assets arrive on site}{The arrival of significant quantities of copper cable, valuable plant, or materials that have been targeted in recent local crime incidents should trigger a review of the security plan, including storage location, CCTV coverage, and access control.}
  \bitem{Following local area threat intelligence}{Any credible intelligence of planned criminal activity at the site, any social media threat directed at the site or its workers, or any Counter Terrorism Policing advisory relating to the area should trigger an immediate security risk assessment review.}
  \bitem{At the beginning of the out-of-hours working period}{Whenever a project moves into an intensive out-of-hours working programme (night shifts, weekend working), the lone worker arrangements must be reviewed to ensure that the check-in system is in operation and that all workers scheduled for out-of-hours attendance have been individually assessed.}
\end{itemize}

%---------------------------- SECTION ----------------------------
\section{Common failure modes}

\begin{itemize}
  \bitem{No check-in system in operation despite a written lone worker policy}{The most common lone worker control failure is the written policy that exists but is not implemented. The check-in system is perceived as bureaucratic overhead and is omitted by managers who judge that the activity is low-risk and that the worker is ``only going to be there for an hour.'' An hour is sufficient for a cardiac arrest, a serious fall, or an intruder assault to become fatal in the absence of timely intervention.}
  \bitem{CCTV installed but not monitored or maintained}{CCTV that is not connected to a monitoring service, is not maintained, or has cameras that have failed without replacement provides a false sense of security. The investigation of a site theft or violent incident invariably reveals that the relevant CCTV camera was not operational or was positioned incorrectly.}
  \bitem{Hoarding inadequate for an urban site with high public footfall}{Standard timber hoarding panels fixed with screw fixings and subject to the accumulated damage of a busy city-centre location become permeable within weeks. Panels knocked by delivery vehicles, panels from which fixings have been removed by persistent intruders, and panels through which children have made access gaps must be identified and repaired within 24 hours.}
  \bitem{Security guard duties unclear and response procedures absent}{Security guards employed on construction sites who have not been briefed on the site hazards (open excavations, unprotected openings, live services areas) and who have no documented procedure for responding to intruders, fire, or lone worker emergencies are both at risk themselves and unable to provide the protection the site requires. The security guard's role must be specified in writing as part of the security plan.}
  \bitem{Violence incidents not reported and not investigated}{Violence at work in construction is significantly under-reported. Operatives and managers who are threatened, pushed, or assaulted on site frequently do not report the incident because they expect it will not be taken seriously, because they fear retaliation, or because they regard it as an occupational norm. Under-reporting prevents the identification of patterns of violence risk and the implementation of controls that might prevent a more serious incident.}
\end{itemize}

\example{Failure Pattern}{A logistics coordinator at a large mixed-use development is working alone in the site offices on a Sunday completing programme data. The site hoarding has a panel that was damaged by a delivery vehicle on Friday and not repaired over the weekend. At 3\,pm an intruder enters through the damaged panel, approaches the site office, and demands the coordinator give him the keys to a telehandler. The coordinator complies and is not physically harmed, but is severely shaken and takes three weeks of stress-related sick leave. The intruder drives the telehandler through the hoarding and abandons it. Investigation reveals: no lone worker risk assessment for Sunday office working; no check-in procedure; hoarding damage identified in the Friday afternoon inspection but repair not actioned; no CCTV coverage of the site access gate closest to the damaged panel. Three Improvement Notices are issued by the HSE for lone worker, CCTV, and hoarding failures.}

%---------------------------- CHAPTER SUMMARY ----------------------------
\newpage
\section{Chapter Summary}

\begin{table}[H]
\centering
\begin{tabularx}{\textwidth}{>{\raggedright\arraybackslash}p{3.2cm}X}
\toprule
\textbf{Assessment Element} & \textbf{Operational Requirement} \\
\midrule
Legal/Professional Basis & Management of Health and Safety at Work Regulations 1999; Health and Safety at Work etc.\ Act 1974; RIDDOR 2013; BS 8484:2022; CDM 2015; Secure by Design; UK GDPR 2018. \\
Method & Five-step lone worker, security, and violence risk assessment; BS 8484 compliant device or app; site security plan; CCTV system; HVM assessment. \\
Controls & Lone worker check-in procedure; BS 8484-compliant monitoring; perimeter fencing to appropriate standard; CCTV coverage of all access points and perimeter; HVM barriers at vehicle access points; violence reporting system; post-incident support. \\
Cadence & At each phase transition; after any security or lone worker incident; when new high-value assets arrive; following local threat intelligence; at commencement of out-of-hours working programmes. \\
Assurance & Monthly security review; CCTV maintenance log; check-in system audit; incident log review; annual lone worker policy and security plan review. \\
\bottomrule
\end{tabularx}
\end{table}

\begin{itemize}
  \bitem{Key takeaway 1}{Lone working in construction is widespread and frequently unrecognised; the site manager working alone before the rest of the site arrives is as much a lone worker as the surveyor in an unoccupied building, and both require a check-in system and an emergency response procedure.}
  \bitem{Key takeaway 2}{BS 8484:2022 compliance for lone worker devices and apps enables emergency response through police grade-1 protocol; this is a significant operational advantage over unmonitored check-in systems and should be specified in any lone worker device procurement.}
  \bitem{Key takeaway 3}{Site security is a construction safety matter as well as a commercial one; the intruder who accesses an active construction site and falls through an unprotected opening, contacts live services, or is struck by plant creates both a criminal liability for the contractor and a moral responsibility for the consequences.}
  \bitem{Key takeaway 4}{Hostile Vehicle Mitigation must be considered on all urban construction sites adjacent to pedestrian areas; temporary concrete or water-filled barriers at vehicle access points are a proportionate and readily available control.}
  \bitem{Key takeaway 5}{Violence at work in construction is under-reported; a reporting culture that treats violence as a normal occupational hazard rather than a reportable incident prevents the identification of risk patterns and the implementation of controls that could prevent the escalation to more serious harm.}
\end{itemize}

\barquote{In Chapter~\ref{chap:CDM Duties, Documentation and Competence Risk Assessment}, we examine the risk assessment implications of the CDM 2015 duty holder framework, key documentation obligations, and competence assurance across the construction project.}
